\section{State-dependent approximate balance}
\label{sec:approx}
Exact task balance can be replaced by a uniform conditional certificate. The reference and the all-strategy effort inequality are left unchanged. The upper program may now have arbitrary history-dependent physical states. Write $\nu_{\theta,h}$ for the virtual report law of an attempt begun at a surviving history $h$, and suppose
\begin{align}
 \|\nu_{\theta,h}-\nu\|_{\TV}&\le a<1,\label{eq:approx-mass}\\
 \left|\E_\theta[q(X)\|Z-T_Jr(X)\|^2\mid h]
       -\int q(x)d_{r(x)}(x)\nu(dx)\right|&\le e.\label{eq:approx-task}
\end{align}
Here total variation between probabilities is $\sup_D|\mu(D)-\nu(D)|$. The second bound is required for every Borel $q:X\to[0,1]$ and every finite-image Borel $r:X\to H$, uniformly in $\theta,h$. It compares the complete task measure, not just report marginals or unweighted conditional variances. The program has $B$ live roles, export multiplicity $D_0$, and conditional expected attempt cost at most $\kappa_+$. The constants are certified before selecting a dictionary.

\begin{theorem}[Approximate acquired-geometry transfer]\label{thm:approx}
Under the preceding conditions, for $\beta\ge1$ with $\beta a<1$, one common program uses at most $B$ live and $D_0m$ terminal roles and satisfies
\begin{equation}\label{eq:approx-basic}
 \sup_\theta\E_\theta C_\tau\le\frac{\kappa_+\beta}{1-\beta a},
 \qquad
 \sup_\theta\E_\theta\|Z-\widehat r\|^2
 \le\frac{\Qm_m(\beta)+\beta e}{1-\beta a}.
\end{equation}
It commits almost surely. Thus the alphabet counts are $B+D_0m$ for state stopping and $\max\{B,D_0m\}$ for event stopping. For a specified $b=N/\kappa_+\ge(1-a)^{-1}$, put $\beta=b/(1+ba)$. At the same physical budget $N$, the upper becomes
\begin{equation}\label{eq:approx-budget}
 (1+ba)\Qm_m\left(\frac b{1+ba}\right)+be.
\end{equation}
For $1\le b<(1-a)^{-1}$ the first-attempt rule instead has cost at most $\kappa_+$ and risk at most $\Qm_m(1)+e$. All converses in Theorem~\ref{thm:main} continue to hold at $b_-=N/\kappa_-$. No independence, renewal, or normalized-filter contraction is assumed.
\end{theorem}
\begin{proof}
Take an attained $(C,h_*)$ at cap $\beta$, and accept with $q=h_*/\beta$. Its reference acceptance mass is $1/\beta$. By \eqref{eq:approx-mass}, its conditional physical acceptance probability on every surviving history is at least $p_-=1/\beta-a>0$. Consequently the probability of reaching attempt $j+1$ is at most $(1-p_-)^j$, by iterated conditioning. In particular commitment is almost sure and the expected number of attempts is at most $p_-^{-1}$. This proof does not condition on successful attempts.

The nearest-center rule is finite-image and measurable. By \eqref{eq:approx-task}, the expected loss contributed by any reached attempt is at most $\Qm_m(\beta)/\beta+e$. Sum the nonnegative terminal contributions over the disjoint commitment events, and sum conditional attempt costs over the surviving attempts. Both sums are bounded by their one-attempt upper times $p_-^{-1}$, proving \eqref{eq:approx-basic}. Lemma~\ref{lem:packing} supplies the alphabet counts. The substitution for $\beta$ has $1-\beta a=(1+ba)^{-1}$ and gives the displayed budget and risk exactly. The restriction on $b$ is precisely $\beta\ge1$. With $q=1$, the first attempt always commits, regardless of the report perturbation; hence the small-budget assertion follows without dividing by an approximate acceptance mass. The converse proof uses only the admitted reference and was not altered.
\end{proof}

\begin{corollary}[A stable geometric regime]\label{cor:approx-geometry}
Suppose the marked geometry of Proposition~\ref{prop:geometry} holds, the effort constants have a fixed positive ratio, and the relevant role widths differ by at most a fixed export factor. If $ba$ is bounded, $b$ is above the feasibility threshold in Theorem~\ref{thm:approx}, and
\[
 be\le C_0\max\{b^{-2/\alpha},(bm)^{-2/(d+\alpha)}\},
\]
then upper and lower risks are comparable to the same marked-geometry profile at $b,m$. These are sufficient stability conditions, not necessary error floors for every experiment.
\end{corollary}
\begin{proof}
The cap $b/(1+ba)$ differs from $b$ by a bounded factor. Apply the two power terms of Proposition~\ref{prop:geometry} to \eqref{eq:approx-budget} and use the unchanged converse and fixed width factors.
\end{proof}

\begin{proposition}[Primitive moment certificates]\label{prop:approx-moments}
For identity exports, let $\mu_0,\mu_1,\mu_2$ be the conditional report, vector first-moment and scalar squared-norm measures of an attempt. Suppose, uniformly over all models and surviving histories,
\[
 |\mu_0-\nu|(X)\le d_0,\quad
 |\mu_1-f\nu|(X)\le d_1,\quad
 |\mu_2-(w+\|f\|^2)\nu|(X)\le d_2,
\]
where $|\cdot|$ denotes total variation of scalar or finite-dimensional vector measures. If $\sup_{r\in H}\|r\|\le R_0$, then \eqref{eq:approx-mass}--\eqref{eq:approx-task} hold with
\[
 a\le d_0/2,\qquad e\le d_2+2R_0d_1+R_0^2d_0.
\]
In particular, a conditional total-variation bound $a$ on the joint report--audit law gives $e=La$ whenever all squared losses are bounded by $L$.
\end{proposition}
\begin{proof}
The two report measures are probabilities, so their probability total variation is half the variation norm of their difference. Expand the loss as $\|Z\|^2-2\langle r(X),Z\rangle+\|r(X)\|^2$ and integrate the three bounded coefficient functions against the indicated signed measures. This proves the first assertion for every $q,r$ simultaneously. The last assertion applies the defining probability total-variation bound to a measurable function taking values in $[0,L]$. Vector coordinates may be taken in the affine hull of $H$; no ambient redundant coordinates are needed.
\end{proof}

\begin{proposition}[Sharp rare-selection scales]\label{prop:sharp-approx}
The factor $b$ multiplying the one-attempt task defect in \eqref{eq:approx-budget} cannot be replaced by a universal smaller-order factor, even with equal report laws and one live role. Likewise the waiting denominator $1-\beta a$ in \eqref{eq:approx-basic} is attained by a zero-loss selection problem. Absence of first-moment balance can prevent every common causal program, not merely a proposed finite mixture, from attaining the reference profile.
\end{proposition}
\begin{proof}
For the first assertion take $H=[-1,1]$, and let every productive call cost one and draw an independent mark uniformly on $[0,1]$. Let $D$ have mass $1/b$. In the reference model the audit is zero on $D$ and is a fresh fair sign off $D$. In a second model it is a fresh fair sign multiplied by $\sqrt{be}$ on $D$ and an ordinary fair sign off $D$, where $0\le be\le1$. Audit coins are unobserved before commitment. Both first moments vanish, all reports agree, and for every acceptance rule and response the difference of unnormalized losses is
\[
 be\int_D q(x)\nu(dx)\le e.
\]
The reference profile is zero, attained by selecting $D$ and predicting zero. In the second model every conditional residual variance is at least $be$, so every almost-surely committing program has risk at least $be$, regardless of its memory. Selecting $D$ commits in expected $b$ calls and attains $be$ uniformly over the two models. Thus the common minimax risk is exactly $be$.

For the denominator take reference report mass $1/\beta$ on a marked region $D$, with zero audit there and a fresh fair sign outside. Perturb its mass to $1/\beta-a>0$ by moving mass outside $D$. This changes the report law by exactly $a$. Every zero-risk rule must commit in $D$, and the reference-count argument at the perturbed independent law forces expected calls at least $(1/\beta-a)^{-1}$. First entry into $D$ attains this value. The case $\beta=1$ can be treated with two report atoms, the second initially null. This establishes sharp waiting amplification, not a universal positive risk penalty for mere report error.

Finally consider three models with exactly the same constant pre-audit reports and no task-changing commands. Their audits are respectively $0,+1,-1$. The reference profile at the first model is zero at every feasible budget. For any common program, the law of its committed output $R$ is the same under the last two models. Their average loss is $1+\E R^2\ge1$. The constant output zero attains worst risk one. This is a complete indistinguishability proof for all legal history-dependent programs. It is a limitation of that particular information structure, not a claim that general causal balance is impossible. The underlying decision-theoretic argument is classical.
\end{proof}
